\section{Le catalogue des trente-cinq}

\subsection{Durées et tailles}

\begin{figure}[H]
\centering
\includegraphics[width=\linewidth]{usdjpy_trades_gantt}
\caption{Chaque trade, du premier au dernier. La longueur de la barre est la
durée de détention, son épaisseur le volume engagé, sa couleur le signe du
résultat. Les résultats supérieurs à 20\,000\,\$ en valeur absolue sont
annotés.}
\end{figure}

La lecture est brutale~: une barre longue tout en haut, puis trente-quatre
barres courtes qui se pressent sur les trois dernières années. La durée médiane
est de \metric{5~jours}, la durée moyenne de \metric{32~jours} --- un écart qui
dit à lui seul que la distribution est dominée par sa queue. Les volumes, eux,
varient peu~: de 7,85 à 24,60~lots, l'essentiel de la variation venant de la
croissance du compte plutôt que d'une modulation du risque.

\subsection{La cascade du résultat}

\begin{figure}[H]
\centering
\includegraphics[width=\linewidth]{usdjpy_trades_waterfall}
\caption{Construction du résultat, dans l'ordre des clôtures. Chaque barre est
un trade~; la ligne sombre suit la balance du sous-compte.}
\end{figure}

La forme n'a rien de commun avec celle de l'or. Ici le résultat est acquis
\emph{d'emblée}~: la troisième clôture, le 14~novembre~2022, fait passer la
balance de 98\,511 à 286\,884~dollars. Tout ce qui suit --- trente-deux trades,
trois ans et demi --- est une oscillation autour de ce niveau, qui se termine à
276\,822\,\$, soit \emph{plus bas} qu'elle n'a commencé. Un investisseur entré
après novembre~2022 n'aurait rien gagné.

\subsection{Durée contre résultat}

\begin{figure}[H]
\centering
\includegraphics[width=\linewidth]{usdjpy_trades_scatter_duration}
\caption{Durée et résultat de chaque trade. La taille du point est le volume
engagé. L'échelle des durées est logarithmique.}
\end{figure}

La relation avec la durée est nette, et de même sens que sur l'or. Les
\metric{quatre} positions tenues plus de cent jours totalisent
$+$333\,928\,\$, soit \metric{189\,\%} du résultat net~; les vingt-deux
positions d'une semaine ou moins perdent ensemble $-$119\,524\,\$. Le moteur ne
gagne pas en tradant, il gagne en tenant --- et il paie cher chacune de ses
tentatives avortées.

\begin{insightbox}
Cette asymétrie a une conséquence opérationnelle directe~: sur ce candidat, ce
n'est pas seulement le gain qui vient des tenues longues, c'est \emph{aussi}
le portage, qui s'accumule jour après jour et ne rémunère que la durée. Les
deux moteurs de résultat pointent dans le même sens, et toute modification qui
raccourcirait les positions --- horizons plus courts, stop plus serré, prise de
profit --- les attaquerait tous les deux à la fois.
\end{insightbox}

\subsection{La distribution}

\begin{figure}[H]
\centering
\includegraphics[width=\linewidth]{usdjpy_trades_distribution}
\caption{Distribution des résultats par trade, gagnants et perdants séparés,
avec la dispersion en encart.}
\end{figure}

Le corps de la distribution est presque symétrique~: dix-huit gagnants,
dix-sept perdants, une médiane à $+$273\,\$, c'est-à-dire pratiquement zéro. La
moyenne, elle, s'établit à $+$5\,052\,\$ par trade. Tout l'écart entre les deux
tient dans le point isolé à droite du graphique --- l'unique trade à
$+$188\,374\,\$. Un lecteur qui jugerait ce candidat sur son trade médian
conclurait qu'il ne fait rien~; un lecteur qui le jugerait sur sa moyenne
extrapolerait un seul événement.

% Généré par scripts/build_usdjpy_trades_figures.py — NE PAS ÉDITER À LA MAIN.
\begin{table}[H]
\centering
\begin{tabular}{@{}lrrrr@{}}
\toprule
Année & Trades & Gagnants & Portage & Net \\
\midrule
2021 & 2 & 1 & +369 & -1,489 \\
2022 & 1 & 1 & +25,863 & +188,374 \\
2023 & 10 & 5 & +20,608 & -26,247 \\
2024 & 8 & 4 & +28,792 & +52,811 \\
2025 & 9 & 3 & +14,698 & -79,320 \\
2026 & 5 & 4 & +21,936 & +42,695 \\
\bottomrule
\end{tabular}
\caption{Résultat par année de clôture}
\end{table}

